\newcounter{stoppagebreaks}
\documentclass{article}

% packages
\usepackage{amsmath, amsthm, thmtools, amsfonts, amssymb, luacode, catchfile, hyperref, ifthen, graphicx, bbm}
\ifcsname c@kobocompile\endcsname
	\usepackage[a5paper, total={1072pt, 1448pt}, margin=10pt, includeheadfoot]{geometry} % set page margins
\else
	\usepackage[a4paper, margin=50pt, includeheadfoot]{geometry}
\fi
\usepackage[shortlabels]{enumitem}
\usepackage[skip=3pt, indent=0pt]{parskip}

% language
\usepackage[bidi=basic, layout=tabular, provide=*]{babel}
\ifcsname c@english\endcsname
	\babelprovide[main, import]{english}
\else
	\babelprovide[main, import, alph=letters.gershayim]{hebrew}
	\setlist[enumerate,2]{label={\alph*.}}
	\babelprovide{rl}
\fi
\babelfont{rm}[Renderer=Harfbuzz]{Libertinus Serif}
\babelfont{sf}{Libertinus Sans}
\babelfont{tt}{Libertinus Mono}

% style
\ifcsname c@stoppagebreaks\endcsname
\else
	\AddToHook{cmd/section/before}{\clearpage}	% Add line break before section
\fi
\linespread{1.3}
\setcounter{secnumdepth}{0}		% Remove default number tags from sections, this won't do well with theorems
\AtBeginDocument{\setlength{\belowdisplayskip}{3pt}}
\AtBeginDocument{\setlength{\abovedisplayskip}{3pt}}
\graphicspath{ {../images/} }

% operators
\DeclareMathOperator\cis{cis}
\DeclareMathOperator\Sp{Span}
\DeclareMathOperator\tr{tr}
\DeclareMathOperator\im{Im}
\DeclareMathOperator\re{Re}
\DeclareMathOperator\diag{diag}
\DeclareMathOperator*\lowlim{\underline{lim}}
\DeclareMathOperator*\uplim{\overline{lim}}
\DeclareMathOperator\rng{rng}
\DeclareMathOperator\Sym{Sym}
\DeclareMathOperator\Arg{Arg}
\DeclareMathOperator\Log{Log}
\DeclareMathOperator\dom{dom}
\DeclareMathOperator\supp{Supp}
\DeclareMathOperator\var{Var}
\DeclareMathOperator\cov{Cov}
\DeclareMathOperator\End{End}
\DeclareMathOperator\id{id}
\DeclareMathOperator\cl{cl}
\DeclareMathOperator\gal{Gal}
\DeclareMathOperator\dist{dist}
\DeclareMathOperator\form{form}
\DeclareMathOperator\term{term}
\DeclareMathOperator\sent{sent}
\DeclareMathOperator\Th{Th}
\DeclareMathOperator\Hom{Hom}
\DeclareMathOperator\Aut{Aut}
\DeclareMathOperator\Char{char}
\DeclareMathOperator\vol{vol}
\DeclareMathOperator\GL{GL}
\DeclareMathOperator\SL{SL}
\DeclareMathOperator\lcm{lcm}

% commands
\newcommand{\cupdot}{\mathbin{\mathaccent\cdot\cup}}
\newcommand{\indicator}{\mathbbm{1}}
\newcommand{\AAA}[0]{\mathbb{A}}
\newcommand{\Aa}[0]{\mathcal{A}}
\newcommand{\BB}[0]{\mathbb{B}}
\newcommand{\Bb}[0]{\mathcal{B}}
\newcommand{\CC}[0]{\mathbb{C}}
\newcommand{\Cc}[0]{\mathcal{C}}
\newcommand{\DD}[0]{\mathbb{D}}
\newcommand{\Dd}[0]{\mathcal{D}}
\newcommand{\EE}[0]{\mathbb{E}}
\newcommand{\FF}[0]{\mathbb{F}}
\newcommand{\Ff}[0]{\mathcal{F}}
\newcommand{\Ii}[0]{\mathcal{I}}
\newcommand{\Gg}[0]{\mathcal{G}}
\newcommand{\HH}[0]{\mathbb{H}}
\newcommand{\Hh}[0]{\mathcal{H}}
\newcommand{\Kk}[0]{\mathcal{K}}
\newcommand{\Ll}[0]{\mathcal{L}}
\newcommand{\Mm}[0]{\mathcal{M}}
\newcommand{\NN}[0]{\mathbb{N}}
\newcommand{\Nn}[0]{\mathcal{N}}
\newcommand{\PP}[0]{\mathbb{P}}
\newcommand{\Pp}[0]{\mathcal{P}}
\newcommand{\QQ}[0]{\mathbb{Q}}
\newcommand{\RR}[0]{\mathbb{R}}
\newcommand{\Rr}[0]{\mathcal{R}}
\newcommand{\Ss}[0]{\mathcal{S}}
\newcommand{\TT}[0]{\mathbb{T}}
\newcommand{\Tt}[0]{\mathcal{T}}
\newcommand{\UU}[0]{\mathbb{U}}
\newcommand{\Uu}[0]{\mathcal{U}}
\newcommand{\Vv}[0]{\mathcal{V}}
\newcommand{\Ww}[0]{\mathcal{W}}
\newcommand{\Xx}[0]{\mathcal{X}}
\newcommand{\Yy}[0]{\mathcal{Y}}
\newcommand{\ZZ}[0]{\mathbb{Z}}
\newcommand{\Zz}[0]{\mathcal{Z}}
\newcommand{\acts}[0]{\circlearrowright}
\newcommand{\explain}[2] {
	\begin{flalign*}
		 && \text{#2} && \text{#1}
	\end{flalign*}
}
\newcommand{\maketitleprint}[1][green]{
	\begin{center}
		\includegraphics[width=6cm]{dragon_#1}
	\end{center}
}

% theorem commands
\newtheoremstyle{c_remark}
	{}	% Space above
	{}	% Space below
	{}% Body font
	{}	% Indent amount
	{\bfseries}	% Theorem head font
	{}	% Punctuation after theorem head
	{.5em}	% Space after theorem head
	{\thmname{#1}\thmnumber{ #2}\thmnote{ \normalfont{\text{(#3)}}}}	% head content
\newtheoremstyle{c_definition}
	{3pt}	% Space above
	{3pt}	% Space below
	{}% Body font
	{}	% Indent amount
	{\bfseries}	% Theorem head font
	{}	% Punctuation after theorem head
	{.5em}	% Space after theorem head
	{\thmname{#1}\thmnumber{ #2}\thmnote{ \normalfont{\text{(#3)}}}}	% head content
\newtheoremstyle{c_plain}
	{3pt}	% Space above
	{3pt}	% Space below
	{\itshape}% Body font
	{}	% Indent amount
	{\bfseries}	% Theorem head font
	{}	% Punctuation after theorem head
	{.5em}	% Space after theorem head
	{\thmname{#1}\thmnumber{ #2}\thmnote{ \text{(#3)}}}	% head content

\ifcsname c@english\endcsname
	\theoremstyle{plain}
	\newtheorem{theorem}{Theorem}[section]
	\newtheorem*{theorem*}{Theorem}
	\newtheorem{lemma}[theorem]{Lemma}
	\newtheorem{proposition}[theorem]{Proposition}
	\newtheorem*{proposition*}{Proposition}

	\theoremstyle{definition}
	\newtheorem{definition}[theorem]{Definition}
	\newtheorem*{definition*}{Definition}
	\newtheorem{example}{Example}[section]
	\newtheorem{exercise}{Exercise}[section]

	\theoremstyle{remark}
	\newtheorem*{remark}{Remark}
	\newtheorem*{solution}{Solution}
	\newtheorem{corollary}[theorem]{Corollary}
	\newtheorem{notation}[theorem]{Notation}
\else
	\theoremstyle{c_plain}
	\newtheorem{theorem}{משפט}[section]
	\newtheorem*{theorem*}{משפט}
	\newtheorem{lemma}[theorem]{למה}
	\newtheorem{proposition}[theorem]{טענה}
	\newtheorem*{proposition*}{טענה}

	\theoremstyle{c_definition}
	\newtheorem{definition}[theorem]{הגדרה}
	\newtheorem*{definition*}{הגדרה}
	\newtheorem{example}{דוגמה}[section]
	\newtheorem{exercise}{תרגיל}[section]

	\theoremstyle{c_remark}
	\newtheorem*{remark}{הערה}
	\newtheorem*{solution}{פתרון}
	\newtheorem{corollary}[theorem]{מסקנה}
	\newtheorem{notation}[theorem]{סימון}
\fi

% Questions related commands
\newcounter{question}
\setcounter{question}{1}
\newcounter{sub_question}
\setcounter{sub_question}{1}

\ifcsname c@english\endcsname
	\newcommand{\question}[1][0]{
		\ifthenelse{#1 = 0}{}{\setcounter{question}{#1}}
		\section{Question \arabic{question}}
		\addtocounter{question}{1}
		\setcounter{sub_question}{1}
	}

	\newcommand{\subquestion}[1][0]{
		\ifthenelse{#1 = 0}{}{\setcounter{sub_question}{#1}}
		\subsection{Part \alph{sub_question}}
		\addtocounter{sub_question}{1}
	}
\else
	\newcommand{\question}[1][0]{
		\ifthenelse{#1 = 0}{}{\setcounter{question}{#1}}
		\section{שאלה \arabic{question}}
		\addtocounter{question}{1}
		\setcounter{sub_question}{1}
	}

	\newcommand{\subquestion}[1][0]{
		\ifthenelse{#1 = 0}{}{\setcounter{sub_question}{#1}}
		\subsection{סעיף \localecounter{letters.gershayim}{sub_question}}
		\addtocounter{sub_question}{1}
	}
\fi

% import lua and start of document
\directlua{common = require ('../common')}

\GetEnv{AUTHOR}

% headers
\author{\AUTHOR}
\date\today

\setlist[enumerate,1]{label={\alph*.}}
% chktex-file 9
% chktex-file 17

\usepackage{pgfplots}
\pgfplotsset{compat=1.18}

\title{מתמטיקה א' --- פתרון תרגיל 6}
\date{}

\begin{document}

\maketitle

\section{שאלה 1}
בדיקת חסימות הפונקציות והוכחת הטענות:
\begin{enumerate}
	\item $f(x) = 5x - 3$ --- אינה חסומה מלמעלה ואינה חסומה מלמטה.
		\begin{itemize}
			\item אינה חסומה מלמעלה: יהי $M \in \RR$. נבחר $x = \frac{M + 4}{5}$, ונקבל:
				\[
					f(x) = 5 \cdot \frac{M + 4}{5} - 3 = M + 4 - 3 = M + 1 > M
				\]
				לכן לא קיים חסם מלמעלה.
			\item אינה חסומה מלמטה: יהי $m \in \RR$. נבחר $x = \frac{m + 2}{5}$, ונקבל:
				\[
					f(x) = 5 \cdot \frac{m + 2}{5} - 3 = m + 2 - 3 = m - 1 < m
				\]
				לכן לא קיים חסם מלמטה.
		\end{itemize}

	\item $f(x) = x^2 + 3$ --- חסומה מלמטה ואינה חסומה מלמעלה.
		\begin{itemize}
			\item חסומה מלמטה: לכל $x \in \RR$ מתקיים $x^2 \ge 0$, ולכן $f(x) = x^2 + 3 \ge 3$. מכאן ש־$3$ הוא חסם מלמטה.
			\item אינה חסומה מלמעלה: יהי $M \in \RR$. נבחר $x = \sqrt{|M|} + 1$, ונקבל:
				\[
					f(x) = {\left(\sqrt{|M|} + 1\right)}^2 + 3 = |M| + 2\sqrt{|M|} + 4 > |M| \ge M
				\]
				לכן הפונקציה אינה חסומה מלמעלה.
		\end{itemize}

	\item $f(x) = \ln(x)$ (בתחום ההגדרה $(0, \infty)$) --- אינה חסומה מלמעלה ואינה חסומה מלמטה.
		\begin{itemize}
			\item אינה חסומה מלמעלה: יהי $M \in \RR$. נבחר $x = e^{|M| + 1} > 0$, ונקבל:
				\[
					f(x) = \ln\left(e^{|M| + 1}\right) = |M| + 1 > M
				\]
				לכן הפונקציה אינה חסומה מלמעלה.
			\item אינה חסומה מלמטה: יהי $m \in \RR$. נבחר $x = e^{-|m| - 1} > 0$, ונקבל:
				\[
					f(x) = \ln\left(e^{-|m| - 1}\right) = -|m| - 1 < -|m| \le m
				\]
				לכן הפונקציה אינה חסומה מלמטה.
		\end{itemize}

	\item $f(x) = 3\sin(x) + 2\cos(x)$ --- חסומה מלמעלה ומלמטה. \\
		לכל $x \in \RR$ מתקיים $|\sin(x)| \le 1$ ו־$|\cos(x)| \le 1$. כל אחד מהערכים אם כך יכול לנוע בין $-2$ ל־$2$ ו־$-3$ ל־$3$, וביחד הם נסכמים לכל היותר ל־$-5$ עד $5$.

	\item $f : [-100, 100] \to \RR, \quad f(x) = e^x + x^7$ --- חסומה מלמעלה ומלמטה. \\
		נראה כי $f$ מונוטונית עולה ממש: יהיו $x_1, x_2 \in [-100, 100]$ כך ש־$x_1 < x_2$. הפונקציה המעריכית $e^x$ עולה ממש ולכן $e^{x_1} < e^{x_2}$, ופונקציית החזקה האי־זוגית $x^7$ עולה ממש ולכן $x_1^7 < x_2^7$. מחיבור אי־השוויונות נקבל:
		\[
			f(x_1) = e^{x_1} + x_1^7 < e^{x_2} + x_2^7 = f(x_2)
		\]
		מכיוון ש־$f$ עולה, לכל $x \in [-100, 100]$ מתקיים $-100 \le x \le 100$, ולכן:
		\[
			f(-100) \le f(x) \le f(100) \implies e^{-100} - 100^7 \le f(x) \le e^{100} + 100^7
		\]
		מכאן ש־$e^{-100} - 100^7$ הוא חסם מלמטה ו־$e^{100} + 100^7$ הוא חסם מלמעלה.
\end{enumerate}

\section{שאלה 2}
נתונות הפונקציות $f(x) = \frac{e^x + e^{-x}}{2}$ ו־$g(x) = \frac{e^x - e^{-x}}{2}$ המוגדרות לכל $x \in \RR$ (תחום סימטרי סביב $0$):
\begin{enumerate}
	\item עבור הפונקציה $f(x) = \frac{e^x + e^{-x}}{2}$:
		\begin{itemize}
			\item זוגיות: לכל $x \in \RR$ מתקיים:
				\[
					f(-x) = \frac{e^{-x} + e^{-(-x)}}{2} = \frac{e^{-x} + e^x}{2} = \frac{e^x + e^{-x}}{2} = f(x)
				\]
				לכן $f$ פונקציה זוגית.
			\item חסומה מלמטה: לכל $x \in \RR$ מתקיים $e^x > 0$ ו־$e^{-x} > 0$, ולכן $f(x) = \frac{e^x + e^{-x}}{2} > 0$. מכאן ש־$0$ הוא חסם מלמטה של $f$.
		\end{itemize}

	\item עבור הפונקציה $g(x) = \frac{e^x - e^{-x}}{2}$:
		\begin{itemize}
			\item אי־זוגיות: לכל $x \in \RR$ מתקיים:
				\[
					g(-x) = \frac{e^{-x} - e^{-(-x)}}{2} = \frac{e^{-x} - e^x}{2} = -\frac{e^x - e^{-x}}{2} = -g(x)
				\]
				לכן $g$ פונקציה אי־זוגית.
			\item אינה חסומה כלל:
				\begin{enumerate}
					\item אינה חסומה מלמעלה: יהי $M \in \RR$. עבור $x > 0$ מתקיים $e^{-x} < 1$, ולכן $g(x) > \frac{e^x - 1}{2}$. נבחר $x = \ln(2|M| + 2) > 0$, ונקבל:
						\[
							g(x) > \frac{e^{\ln(2|M| + 2)} - 1}{2} = \frac{2|M| + 2 - 1}{2} = |M| + \frac{1}{2} > M
						\]
					\item אינה חסומה מלמטה: יהי $m \in \RR$. מכיוון ש־$g$ אינה חסומה מלמעלה, קיים $x_0 \in \RR$ כך ש־$g(x_0) > -m$. מאי־זוגיות הפונקציה נקבל עבור $x = -x_0$:
						\[
							g(-x_0) = -g(x_0) < -(-m) = m
						\]
				\end{enumerate}
				לכן $g$ אינה חסומה מלמעלה ואינה חסומה מלמטה.
		\end{itemize}
\end{enumerate}

\section{שאלה 3}
פתרון המשוואות הטריגונומטריות (בכל הסעיפים $k \in \mathbb{Z}$):
\begin{enumerate}
	\item $\sin x = \frac{1}{2}$: \\
		הפתרונות הבסיסיים במחזור $[0, 2\pi)$ הם $x = \frac{\pi}{6}$ ו־$x = \pi - \frac{\pi}{6} = \frac{5\pi}{6}$, ולכן הפתרון הכללי הוא:
		\[
			x = \frac{\pi}{6} + 2\pi k \quad \text{או} \quad x = \frac{5\pi}{6} + 2\pi k \quad (k \in \mathbb{Z})
		\]

	\item $\cos(2x) = \frac{1}{2}$: \\
		מפתרון משוואת הקוסינוס הבסיסית נקבל:
		\[
			2x = \pm \frac{\pi}{3} + 2\pi k \quad (k \in \mathbb{Z})
		\]
		נחלק את שני האגפים ב־$2$:
		\[
			x = \pm \frac{\pi}{6} + \pi k \quad (k \in \mathbb{Z})
		\]

	\item $\tan x = 0$: \\
		הפונקציה $\tan x = \frac{\sin x}{\cos x}$ מוגדרת כאשר $\cos x \ne 0$ ומתאפסת כאשר $\sin x = 0$ (ובנקודות אלו $\cos x = \pm 1 \ne 0$). לכן הפתרון הכללי הוא:
		\[
			x = \pi k \quad (k \in \mathbb{Z})
		\]
\end{enumerate}

\section{שאלה 4}
הוכחת מחזוריות ומציאת גודל המחזור $T > 0$:
\begin{enumerate}
	\item $f(x) = \sin(\pi x)$ --- גודל המחזור הוא $T = 2$. \\
		המחזור היסודי של $\sin(u)$ הוא $2\pi$, ולכן נדרוש $\pi T = 2\pi \implies T = 2$.
		אכן, לכל $x \in \RR$ מתקיים:
		\[
			f(x + 2) = \sin\big(\pi(x + 2)\big) = \sin(\pi x + 2\pi) = \sin(\pi x) = f(x)
		\]

	\item $f(x) = \tan\left(21\frac{1}{3}x - 15\right)$ --- גודל המחזור הוא $T = \frac{3\pi}{64}$. \\
		נרשום $21\frac{1}{3} = \frac{64}{3}$, כלומר $f(x) = \tan\left(\frac{64}{3}x - 15\right)$. המחזור היסודי של $\tan(u)$ הוא $\pi$, ולכן נדרוש $\frac{64}{3}T = \pi \implies T = \frac{3\pi}{64}$.
		אכן, לכל $x$ בתחום ההגדרה של $f$ (תחום הנשמר תחת הזזה ב־$\frac{3\pi}{64}$) מתקיים:
		\[
			f\left(x + \frac{3\pi}{64}\right) = \tan\left(\frac{64}{3}\left(x + \frac{3\pi}{64}\right) - 15\right) = \tan\left(\left(\frac{64}{3}x - 15\right) + \pi\right) = \tan\left(\frac{64}{3}x - 15\right) = f(x)
		\]

	\item $f(x) = \tan(3x) - \sin(5x)$ --- גודל המחזור הוא $T = 2\pi$. \\
		המחזור היסודי של $\tan(3x)$ הוא $T_1 = \frac{\pi}{3}$, והמחזור היסודי של $\sin(5x)$ הוא $T_2 = \frac{2\pi}{5}$. המחזור המשותף המינימלי הוא הכפולה המשותפת המינימלית החיובית של $T_1$ ו־$T_2$:
		\[
			n \cdot \frac{\pi}{3} = m \cdot \frac{2\pi}{5} \implies 5n = 6m
		\]
		הפתרון המינימלי במספרים טבעיים הוא $n = 6, m = 5$, ולכן $T = 6 \cdot \frac{\pi}{3} = 2\pi$.
		אכן, לכל $x$ בתחום ההגדרה של $f$ מתקיים:
		\[
			f(x + 2\pi) = \tan\big(3(x + 2\pi)\big) - \sin\big(5(x + 2\pi)\big) = \tan(3x + 6\pi) - \sin(5x + 10\pi) = \tan(3x) - \sin(5x) = f(x)
		\]
\end{enumerate}

\section{שאלה 5}
גרפים המתארים את ההרחבות המחזוריות:
\begin{enumerate}
	\item $f : [0, 10) \to \RR, \quad f(x) = x^2, \quad L = 10$: \\
		בכל קטע $[10k, 10k + 10)$ עבור $k \in \mathbb{Z}$ הגרף הוא העתק של הפרבולה מ־$(10k, 0)$ עד $(10k + 10, 100)$.
		\begin{otherlanguage}{english}
			\begin{center}
				\begin{tikzpicture}[xscale=0.35, yscale=0.032]
					\draw[->] (-11.5, 0) -- (22, 0) node[right] {$x$};
					\draw[->] (0, -8) -- (0, 115) node[above] {$y$};
					\draw[dashed, gray] (-10, 100) -- (20, 100);
					\node[left] at (0, 100) {$100$};
					\foreach \x in {-10, 10, 20} {
						\draw (\x, -2.5) -- (\x, 2.5);
						\node[below] at (\x, -2.5) {$\x$};
					}

					\draw[thick, blue, domain=-10:0, samples=50] plot (\x, {(\x+10)^2});
					\node[circle, fill=blue, inner sep=1.5pt] at (-10, 0) {};
					\node[circle, draw=blue, fill=white, inner sep=1.5pt] at (0, 100) {};

					\draw[thick, blue, domain=0:10, samples=50] plot (\x, {(\x)^2});
					\node[circle, fill=blue, inner sep=1.5pt] at (0, 0) {};
					\node[circle, draw=blue, fill=white, inner sep=1.5pt] at (10, 100) {};

					\draw[thick, blue, domain=10:20, samples=50] plot (\x, {(\x-10)^2});
					\node[circle, fill=blue, inner sep=1.5pt] at (10, 0) {};
					\node[circle, draw=blue, fill=white, inner sep=1.5pt] at (20, 100) {};
				\end{tikzpicture}
			\end{center}
		\end{otherlanguage}

	\item $f : [-5, 5] \to \RR, \quad f(x) = x^2, \quad L = 10$:
		\begin{otherlanguage}{english}
			\begin{center}
				\begin{tikzpicture}[xscale=0.35, yscale=0.11]
					\draw[->] (-16.5, 0) -- (17.5, 0) node[right] {$x$};
					\draw[->] (0, -2.5) -- (0, 30) node[above] {$y$};
					\draw[dashed, gray] (-15, 25) -- (15, 25);
					\node[above left] at (0, 25) {$25$};
					\foreach \x in {-15, -5, 5, 15} {
						\draw (\x, -0.8) -- (\x, 0.8);
						\node[below] at (\x, -0.8) {$\x$};
					}

					\draw[thick, blue, domain=-15:-5, samples=50] plot (\x, {(\x+10)^2});

					\draw[thick, blue, domain=-5:5, samples=50] plot (\x, {(\x)^2});

					\draw[thick, blue, domain=5:15, samples=50] plot (\x, {(\x-10)^2});
				\end{tikzpicture}
			\end{center}
		\end{otherlanguage}

	\item $f : [-1, 2) \to \RR, \quad f(x) = |x|, \quad L = 3$: \\
		בכל קטע $[3k - 1, 3k + 2)$ הגרף מתחיל בנקודה ב־$(3k - 1, 1)$, יורד ל־$(3k, 0)$ ועולה עד לנקודה ב־$(3k + 2, 2)$.
		\begin{otherlanguage}{english}
			\begin{center}
				\begin{tikzpicture}[xscale=1.05, yscale=1.05]
					\draw[->] (-4.6, 0) -- (5.8, 0) node[right] {$x$};
					\draw[->] (0, -0.3) -- (0, 2.5) node[above] {$y$};
					\draw[dashed, gray] (-4, 2) -- (5, 2);
					\node[left] at (0, 1) {$1$};
					\node[above left] at (0, 2) {$2$};
					\foreach \x in {-4, -3, -1, 2, 3, 5} {
						\draw (\x, -0.08) -- (\x, 0.08);
						\node[below] at (\x, -0.08) {$\x$};
					}

					\draw[thick, blue] (-4, 1) -- (-3, 0) -- (-1, 2);
					\node[circle, fill=blue, inner sep=1.5pt] at (-4, 1) {};
					\node[circle, draw=blue, fill=white, inner sep=1.5pt] at (-1, 2) {};

					\draw[thick, blue] (-1, 1) -- (0, 0) -- (2, 2);
					\node[circle, fill=blue, inner sep=1.5pt] at (-1, 1) {};
					\node[circle, draw=blue, fill=white, inner sep=1.5pt] at (2, 2) {};

					\draw[thick, blue] (2, 1) -- (3, 0) -- (5, 2);
					\node[circle, fill=blue, inner sep=1.5pt] at (2, 1) {};
					\node[circle, draw=blue, fill=white, inner sep=1.5pt] at (5, 2) {};
				\end{tikzpicture}
			\end{center}
		\end{otherlanguage}
\end{enumerate}

\end{document}
